\section{Immunological control}

\subsection{Control formulation}
Represent control via $\mu$ (variational switch) and an adaptive field
\[
u(\mathbf{x},t)=\min\!\left(u_{\max},\,u_0\, \frac{c}{i+\varepsilon}\right)
\]
to prioritize high-$c$, low-$i$ regions. Coupling can enhance $\beta$ (cytotoxicity) or enter as a source in the immune equation.

\subsection{Control strategies}
\begin{itemize}
    \item Constant vs adaptive dosing.
    \item Spatial targeting: stronger control where $c$ is high and $i$ is low.
    \item Timing protocols: pulsed vs continuous application.
\end{itemize}

\subsection{Efficiency and cost}
Define a cost functional (drug load, $\int u\, d\mathbf{x}\,dt$) vs benefit (peak reduction, time to extinction, total tumor burden $\int c\,dt$, correlation length of $c$). Compare $\mu=0,1$ and adaptive $u$; report Pareto curves cost vs suppression. Qualitatively: $\mu=1$ reduces fragmentation and peak burden; adaptive $u$ focuses resources on high-$c$, low-$i$ zones, improving containment for similar cost.

\setcounter{table}{1} % ensure Table II numbering under revtex onecolumn
\begin{table*}[!htbp]
\centering
\begin{tabular}{lccc}
\toprule
Protocol & Peak $c$ reduction & Time to extinction & Cost (relative) \\
\midrule
$\mu=0$ (no control) & baseline & -- & 0 \\
$\mu=1$ (binary) & $\downarrow$ moderate & shorter vs baseline & medium \\
Adaptive $u$ (weak Allee) & $\downarrow$ strong & shortest & medium--high \\
Adaptive $u$ (strong Allee) & $\downarrow$ strong & shortest & medium--high \\
\bottomrule
\end{tabular}
\caption{Table II (qualitative): control efficiency vs cost. Replace with measured metrics (peak, extinction time, integrated $u$) when available.}
\label{tab:control-cost}
\end{table*}

\begin{figure}[!htbp]
    \centering
\includegraphics[draft=false,width=0.6\linewidth]{estabilidad_lineal_strong_mu1.png}
    \caption{Ejemplo de perfil espectral bajo control ($\mu=1$) para Allee fuerte: se atenúan modos de alta frecuencia, reduciendo fragmentación.}
\end{figure}

% Fig. 6: control field $u(x,t)$ applied on $c$-rich, $i$-poor regions.

% Fig. 6: control field $u(x,t)$. Tabla II: coste vs supresión.

